% =====================================================================
%  ПРЕЗЕНТАЦИЯ ПО ГЛАВЕ 2
%  «Векторные данные в Python: OGR, шейп-файлы и расстояние по большому кругу»
%
%  Материал: Python Geospatial Development, 3rd Edition (Erik Westra),
%            репозиторий, папка Chapter 02 — пять файлов
%            example01.py … example05.py.
%
%  ВСЕ ЧИСЛА, ИМЕНА ОБЪЕКТОВ И ВЫВОДЫ ПРОГРАММ В ПРЕЗЕНТАЦИИ РЕАЛЬНЫЕ:
%  пять примеров запущены на настоящем шейп-файле
%  tl_2014_us_state.shp (US Census Bureau, TIGER/Line 2014).
%
%  ---------------------------------------------------------------
%  КАК КОМПИЛИРОВАТЬ В TeXStudio
%  ---------------------------------------------------------------
%  Нужен движок с поддержкой Unicode-шрифтов:
%     Options -> Configure TeX -> Compiler -> выбрать LuaLaTeX
%     (подойдёт также XeLaTeX; pdflatex НЕ подходит — нет кириллицы).
%     Затем Tools -> Compile. Сборка делается в два прохода
%     (для оглавления/навигации), TeXStudio сделает это сам.
%
%  Картинки лежат в подпапке figures/ рядом с этим файлом. Если какой-то
%  из них потерялся, на его месте появится рамка-заглушка — сборка не падает.
% =====================================================================
\documentclass[aspectratio=169,10pt]{beamer}

% ---------- шрифты и язык ----------
\usepackage{fontspec}
\usepackage{polyglossia}
\setdefaultlanguage{russian}
\setmainfont{DejaVu Serif}
\setsansfont{DejaVu Sans}
\setmonofont{DejaVu Sans Mono}

% ---------- оформление ----------
\usepackage{tikz}
\usetikzlibrary{arrows.meta,positioning,calc,shapes.geometric,fit,backgrounds}
\usepackage{booktabs}
\usepackage{array}
\usepackage{graphicx}
\usepackage{fvextra}
\usepackage{caption}
\usepackage{hyperref}
\hypersetup{colorlinks=true,linkcolor=blue!60!black,urlcolor=blue!60!black}

\usetheme{default}
\usecolortheme{default}
\setbeamercolor{structure}{fg=blue!55!black}
\setbeamercolor{frametitle}{fg=white,bg=blue!55!black}
\setbeamercolor{block title}{fg=white,bg=blue!50!black}
\setbeamercolor{block body}{fg=black,bg=blue!5}
\setbeamertemplate{navigation symbols}{}
\setbeamertemplate{itemize item}{\color{blue!60!black}$\blacktriangleright$}
\setbeamertemplate{itemize subitem}{\color{blue!45!black}$\bullet$}
\setbeamerfont{frametitle}{size=\large,series=\bfseries}
\setbeamersize{text margin left=6mm,text margin right=6mm}
\setlength{\parskip}{2pt}
\setbeamercovered{transparent}

% ---------- блоки кода ----------
\fvset{
  fontsize=\footnotesize,
  baselinestretch=0.92,
  frame=single,
  rulecolor=\color{gray!45},
  fillcolor=\color{gray!4},
  numbers=left,
  numbersep=7pt,
  firstnumber=1,
  stepnumber=1,
  numberblanklines=true,
  breaklines=true,
  breaksymbolleft=\textcolor{red!60}{\tiny$\hookrightarrow$},
  breakanywhere=true,
  tabsize=4,
  xleftmargin=13pt,
  framesep=3pt,
  showspaces=false
}
\makeatletter
\renewcommand{\theFancyVerbLine}{\textcolor{gray!65}{\scriptsize\arabic{FancyVerbLine}}}
\makeatother

\definecolor{codegreen}{rgb}{0.10,0.45,0.20}
\definecolor{codeblue}{rgb}{0.05,0.30,0.60}
\definecolor{codered}{rgb}{0.70,0.15,0.10}
\definecolor{numblue}{rgb}{0.10,0.25,0.60}

\newcommand{\hl}[1]{\textcolor{codeblue}{\textbf{#1}}}
\newcommand{\cmd}[1]{\textcolor{codegreen}{\textbf{#1}}}
\newcommand{\num}[1]{\textcolor{codered}{\textbf{#1}}}
\newcommand{\out}[1]{\textcolor{black!80}{\scriptsize\texttt{#1}}}

% ---------- картинки ----------
% Проверка существования файла выполняется ОДИН РАЗ в преамбуле:
% \IfFileExists нельзя вызывать прямо внутри кадра beamer ---
% он ломает внутренний просмотр \end{frame}.
\makeatletter
\newcommand{\picdef}[3]{%
  \expandafter\gdef\csname pic@#1\endcsname{%
    \IfFileExists{#2}{\includegraphics[#3]{#2}}%
      {\fbox{\parbox[c][4.2cm][c]{0.92\linewidth}{\centering\footnotesize
          Файл \texttt{#2} не найден.\\[0.3em]
          Он должен лежать в папке \texttt{figures/} рядом с этим файлом.}}}}%
}
\makeatother
\picdef{us}{figures/fig_us_states.png}{width=0.56\linewidth}
\picdef{nm}{figures/fig_nm_polygon.png}{width=0.55\linewidth}
\picdef{az}{figures/fig_az_extremes.png}{width=0.29\linewidth}
\picdef{gc}{figures/fig_great_circle.png}{width=0.52\linewidth}
\picdef{dc}{figures/fig_distance_compare.png}{width=0.80\linewidth}
\newcommand{\usePic}[1]{\expandafter\csname pic@#1\endcsname}

% ---------- титульный слайд ----------
\title[Глава 2. Векторные данные в Python]{\textbf{Векторные данные в Python}}
\subtitle{OGR, шейп-файлы, разбор геометрии\\ и расстояние по большому кругу}
\author{Разбор материалов главы 2}
\date{Репозиторий \texttt{Python-Geospatial-Development-Third-Edition}\\
      \small папка \texttt{Chapter 02}: \texttt{example01.py} \ldots\ \texttt{example05.py}\\
      \small Продолжительность занятия: 90 минут}

\begin{document}

% ======================= 1. ТИТУЛ =======================
\begin{frame}[plain]
  \vspace{0.3cm}
  \begin{center}
    {\LARGE\bfseries \inserttitle\par}
    \vspace{0.2cm}
    {\large \insertsubtitle\par}
  \end{center}
  \vskip0.2cm
  \begin{center}
    \usePic{us}\\[-0.1cm]
    {\footnotesize Шейп-файл \texttt{tl\_2014\_us\_state.shp}: 56 объектов,
      855\,679 вершин, NAD83 (EPSG:4269)}
  \end{center}
  \vfill
  \begin{center}\small\insertauthor\\[0.3em]\insertdate\end{center}
\end{frame}

% ======================= 2. ПЛАН =======================
\begin{frame}[fragile]{План занятия (90 минут)}
\scriptsize
\renewcommand{\arraystretch}{0.92}
\begin{columns}[T,onlytextwidth]
  \begin{column}{0.5\textwidth}
    \textbf{Часть I. Данные и формат}\\[0.1em]
    \begin{tabular}{@{}r@{\quad}p{0.78\linewidth}@{}}
      \toprule
      Время & Тема \\
      \midrule
      5 мин & Что лежит в папке \texttt{Chapter 02} \\
      5 мин & Шейп-файл: пять файлов, бинарный формат \\
      4 мин & СК, атрибуты, объём данных \\
      \midrule
      \multicolumn{2}{@{}l}{\textbf{Часть II. Чтение через OGR}} \\
      \midrule
      6 мин & Объектная модель OGR \\
      8 мин & \texttt{example01}: слои, СК, объекты \\
      7 мин & \texttt{example02}: атрибуты и геометрия \\
      8 мин & \texttt{example03}: рекурсивный разбор \\
      7 мин & \texttt{example04}: крайние точки, производительность \\
    \end{tabular}
  \end{column}
  \begin{column}{0.48\textwidth}
    \textbf{Часть III. Геодезия}\\[0.1em]
    \begin{tabular}{@{}r@{\quad}p{0.74\linewidth}@{}}
      \toprule
      Время & Тема \\
      \midrule
      2 мин & Зачем нам расстояние \\
      9 мин & \texttt{example05}: формула гаверсинуса \\
      6 мин & Геометрия формулы и её вывод \\
      7 мин & Проверка на реальных данных \\
      \bottomrule
    \end{tabular}

    \vspace{0.4em}
    \begin{block}{Итоги занятия}\scriptsize
      \begin{itemize}\itemsep0pt
        \item объектная модель OGR и рекурсия по геометрии;
        \item устройство шейп-файла «изнутри»;
        \item геодезия на сфере и её границы.
      \end{itemize}
    \end{block}
  \end{column}
\end{columns}
\end{frame}

% ======================= 3. ЧТО В ПАПКЕ =======================
\begin{frame}[fragile]{Что лежит в папке \texttt{Chapter 02}}
\footnotesize
В папке пять коротких программ и \emph{ни одного файла данных}.

\vspace{0.2em}
\begin{tabular}{@{}p{0.24\linewidth}p{0.72\linewidth}@{}}
  \toprule
  Файл & Что показывает \\
  \midrule
  \texttt{example01.py} & Обход источника данных: слои, система координат,
    количество объектов и их имена \\
  \texttt{example02.py} & Атрибуты одного объекта и имя его геометрии \\
  \texttt{example03.py} & Рекурсивный «разбор» геометрии: дерево типов
    и числа точек \\
  \texttt{example04.py} & Рекурсивный поиск самой северной и южной
    вершины объекта \\
  \texttt{example05.py} & Расстояние между двумя точками по большому кругу
    (формула гаверсинуса) \\
  \bottomrule
\end{tabular}

\vspace{0.6em}
\begin{columns}[T,onlytextwidth]
  \begin{column}{0.48\textwidth}
    \begin{block}{Данные, которые нужны примерам}\small
      Шейп-файл \texttt{tl\_2014\_us\_state.shp} --- границы штатов США
      из набора \emph{TIGER/Line 2014} (US Census Bureau).
      \medskip

      \textbf{Все числа в этой презентации} получены запуском этих пяти
      файлов именно на этом шейп-файле.
    \end{block}
  \end{column}
  \begin{column}{0.5\textwidth}
    {\scriptsize Все примеры запускаются «как есть», если рядом с ними
     лежит шейп-файл \texttt{tl\_2014\_us\_state.shp}.}

    \vspace{0.3em}
    \begin{alertblock}{Главная мысль}\small
      Пять примеров образуют один сюжет: \textbf{открыть файл $\to$
      понять, что внутри $\to$ рекурсивно обойти геометрию $\to$
      посчитать по ней географию}.
    \end{alertblock}
  \end{column}
\end{columns}
\end{frame}

% ======================= 4. РАЗДЕЛ I =======================
\begin{frame}
  \begin{center}
    {\LARGE\bfseries Часть I. Данные и формат}\\[0.5em]
    {\large Шейп-файл: что это за файл и как он устроен}\\[1.5em]
    {\small \textcolor{gray}{~25 минут}}
  \end{center}
\end{frame}

% ======================= 5. ОБЗОР ДАННЫХ =======================
\begin{frame}{Данные: границы 56 единиц США}
\centering
\usePic{us}\\[0.1cm]
\footnotesize
На карте выделены два объекта, с которыми работают примеры:
\textcolor{codeblue}{\textbf{feature 12 --- New Mexico}} и
\textcolor{codeblue}{\textbf{feature 55 --- Arizona}}.
Верхний правый фрагмент иллюстрирует первую проблему главы:
\textcolor{numblue}{Alaska} --- это \emph{MULTIPOLYGON} из 50 отдельных
полигонов, а \textcolor{numblue}{Hawaii} --- из 17.
\end{frame}

% ======================= 6. ПЯТЬ ФАЙЛОВ =======================
\begin{frame}[fragile]{Шейп-файл = не один файл, а пять}
\footnotesize
\begin{columns}[T,onlytextwidth]
  \begin{column}{0.54\textwidth}
    \captionof{table}{\scriptsize Состав набора на диске (реальные размеры)}
    \begin{tabular}{@{}lrl@{}}
      \toprule
      Файл & Размер & Назначение \\
      \midrule
      \texttt{.shp} & 13\,694\,480 & геометрии \\
      \texttt{.shx} & 548         & индекс записей \\
      \texttt{.dbf} & 10\,450     & атрибуты \\
      \texttt{.prj} & 165         & система координат (WKT) \\
      \texttt{...}  &             & также \texttt{.cpg}, \texttt{.qpj} \\
      \bottomrule
    \end{tabular}
    \vspace{0.6em}
    \begin{itemize}\itemsep1pt
      \item \textbf{геометрия и атрибуты разнесены} по разным файлам;
      \item \texttt{.shx} --- индекс: по номеру объекта сразу находится
            его смещение в \texttt{.shp};
      \item это главный аргумент «за» OGR: формат разбирает библиотека,
            а не вы.
    \end{itemize}
  \end{column}
  \begin{column}{0.44\textwidth}
    \begin{block}{А если файла нет?}
      \small
      \texttt{example01.py} вызовет
      \texttt{ogr.Open("tl\_2014\_us\_state.shp")} и упадёт с \texttt{NULL}
      в переменной \texttt{shapefile}.
    \end{block}
    \vspace{0.2em}
    {\scriptsize Проверка, которую стоит писать в каждом проекте:}
    \begin{Verbatim}[numbers=none,fontsize=\scriptsize]
from osgeo import ogr, gdal
gdal.UseExceptions()      # GDAL >= 3.7

ds = ogr.Open("tl_2014_us_state.shp")
if ds is None:
    raise IOError("файл не найден")
\end{Verbatim}
  \end{column}
\end{columns}
\end{frame}

% ======================= 7. ВНУТРЕННЕЕ УСТРОЙСТВО =======================
\begin{frame}{Что внутри \texttt{.shp}: заголовок и записи}
\footnotesize
\begin{columns}[T,onlytextwidth]
  \begin{column}{0.51\textwidth}
    \captionof{table}{\scriptsize Заголовок \texttt{.shp} --- первые 100 байт}
    \begin{tabular}{@{}p{0.28\linewidth}p{0.68\linewidth}@{}}
      \toprule
      Смещение & Значение в нашем файле \\
      \midrule
      0--3   & File code = \texttt{9994} (главный файл) \\
      24--27 & Длина файла: 6\,847\,240 16-битных слов \\
      28--31 & Version = 1000 \\
      32--35 & Shape type = \textbf{5} (PolygonZ) \\
      36--99 & Extent: $-179.231086,\;179.859681,$\\
             & $\phantom{Extent:}$-14.601813,\;71.441059 \\
      \bottomrule
    \end{tabular}
    \vspace{0.4em}
    {\scriptsize Дальше --- \textbf{56 записей}: номер записи (8 байт),
    длина содержимого (8 байт) и координаты, по 16 байт на точку
    (тип \texttt{double}).}
  \end{column}
  \begin{column}{0.47\textwidth}
    \captionof{table}{\scriptsize \texttt{.shx}: $100 + 56\times8 = 548$ байт}
    \begin{tabular}{@{}p{0.34\linewidth}p{0.6\linewidth}@{}}
      \toprule
      Запись & Смещение и длина \\
      \midrule
      № 0  & 50 слов = 100 байт; \\
            & 281\,280 слов = 562\,560 байт \\
      № 55 & 6\,768\,412 слов = 13\,536\,824 байта \\
      \bottomrule
    \end{tabular}
    \vspace{0.5em}
    \begin{block}{Как читается \texttt{GetFeature(12)}}\small
      Из \texttt{.shx} берутся 8 байт с позиции $100+12\times8$.
      Первое число, умноженное на 2, --- смещение записи в \texttt{.shp};
      второе --- длина её содержимого.
    \end{block}
    \vspace{0.2em}
    {\scriptsize \textbf{Вывод:} порядок объектов фиксирован и
     \emph{не} алфавитный.}
  \end{column}
\end{columns}
\end{frame}

% ======================= 8. ТАБЛИЦА АТРИБУТОВ =======================
\begin{frame}{Что лежит в атрибутах: 14 полей DBF}
\footnotesize
\captionof{table}{\scriptsize Схема слоя \texttt{tl\_2014\_us\_state} (полный состав полей)}
\begin{tabular}{@{}llllll@{}}
  \toprule
  Поле & Тип & Поле & Тип & Поле & Тип \\
  \midrule
  REGION   & String(2)    & GEOID    & String(2)    & ALAND    & Integer64 \\
  DIVISION & String(2)    & STUSPS   & String(2)    & AWATER   & Integer64 \\
  STATEFP  & String(2)    & NAME     & String(100)  & INTPTLAT & String(11) \\
  STATENS  & String(8)    & LSAD     & String(2)    & INTPTLON & String(12) \\
  MTFCC    & String(5)    & FUNCSTAT & String(1)    &          & \\
  \bottomrule
\end{tabular}

\vspace{0.6em}
\begin{columns}[T,onlytextwidth]
  \begin{column}{0.52\textwidth}
    \begin{itemize}\itemsep1pt
      \item \textbf{STATEFP / GEOID} --- двухзначный код штата по FIPS;
      \item \textbf{NAME} --- название: \texttt{"New Mexico"};
      \item \textbf{ALAND / AWATER} --- площадь суши и воды в м$^2$:
            у Нью-Мэксико $314\,161\,410\,324$ м$^2$ суши;
      \item \textbf{INTPTLAT / INTPTLON} --- «внутренняя точка»,
            хранится \emph{строкой}, потому что содержит точку и знак.
    \end{itemize}
  \end{column}
  \begin{column}{0.46\textwidth}
    \begin{alertblock}{Подводный камень}\small
      В DBF нет типов «строка с точкой» и «число с плавающей точкой»:
      \texttt{ALAND} --- это \texttt{Integer64}, а \texttt{INTPTLAT} ---
      \texttt{String}. Преобразовывать придётся вручную:
      \texttt{float(f.GetField("INTPTLAT"))}.
    \end{alertblock}
  \end{column}
\end{columns}
\end{frame}

% ======================= 9. СИСТЕМА КООРДИНАТ =======================
\begin{frame}{Система координат: NAD83 (EPSG:4269)}
\footnotesize
\begin{columns}[T,onlytextwidth]
  \begin{column}{0.5\textwidth}
    \begin{block}{Что сообщает программа}\small
      \out{Layer 0 has spatial reference}\\
      \out{+proj=longlat +datum=NAD83 +no\_defs}
    \end{block}
    \vspace{0.2em}
    \begin{itemize}\itemsep1pt
      \item Координаты в файле --- \textbf{градусы}, а \emph{не} метры;
      \item \texttt{longlat} --- географическая (угловая) СК;
      \item \texttt{datum=NAD83} --- эллипсоид и система отсчёта:
            \textbf{North American Datum 1983}, код \textbf{EPSG:4269};
      \item \texttt{+no\_defs} --- «дополнительных определений нет».
    \end{itemize}
  \end{column}
  \begin{column}{0.48\textwidth}
    \captionof{table}{\scriptsize Что ещё лежит в объекте СК}
    \begin{tabular}{@{}p{0.5\linewidth}p{0.42\linewidth}@{}}
      \toprule
      Вызов & Результат \\
      \midrule
      \texttt{GetName()}        & \texttt{NAD83} \\
      \texttt{GetAuthorityName} & \texttt{EPSG} \\
      \texttt{GetAuthorityCode} & \texttt{4269} \\
      \texttt{ExportToProj4()}  & строка PROJ.4 \\
      \texttt{ExportToWkt()}   & полное WKT-описание \\
      \bottomrule
    \end{tabular}
    \vspace{0.5em}
    \begin{alertblock}{Зачем это знать}\small
      Расстояние в градусах --- не метры. Отсюда растёт вся часть III:
      чтобы получить \textbf{километры}, формула должна знать радиус Земли.
    \end{alertblock}
  \end{column}
\end{columns}
\end{frame}

% ======================= 10. ОБЪЕКТНАЯ МОДЕЛЬ =======================
\begin{frame}{Объектная модель OGR}
\centering
\resizebox{\linewidth}{!}{%
\begin{tikzpicture}[
  font=\scriptsize,
  box/.style={draw,rounded corners=2pt,minimum height=8mm,align=center,
              fill=blue!4,inner sep=4pt,text width=28mm},
  arr/.style={-{Stealth[length=2mm]},thick,draw=blue!60!black}
]
  \node[box] (ds) {\textbf{DataSource}\\[-1mm]\texttt{ogr.Open(...)}\\[-1mm]
                 «источник данных»};
  \node[box,right=11mm of ds] (lyr) {\textbf{Layer}\\[-1mm]\texttt{GetLayer(0)}\\[-1mm]
                 «слой»};
  \node[box,right=11mm of lyr] (ftr) {\textbf{Feature}\\[-1mm]\texttt{GetFeature(12)}\\[-1mm]
                 «объект / запись»};
  \node[box,right=11mm of ftr,fill=orange!12,draw=orange!70!black] (geo)
                 {\textbf{Geometry}\\[-1mm]\texttt{GetGeometryRef()}\\[-1mm]
                 «геометрия»};

  \draw[arr] (ds) -- (lyr);
  \draw[arr] (lyr) -- (ftr);
  \draw[arr] (ftr) -- (geo);

  \node[below=8mm of ds,align=center,text width=30mm,font=\scriptsize]
        {число слоёв: \texttt{GetLayerCount()}\\пространственный\\фильтр (опц.)};
  \node[below=8mm of lyr,align=center,text width=30mm,font=\scriptsize]
        {имя слоя, СК, число\\объектов, границы,\\предел \texttt{SetFeature()}};
  \node[below=8mm of ftr,align=center,text width=30mm,font=\scriptsize]
        {поля атрибутов,\\FID, геометрия};
  \node[below=8mm of geo,align=center,text width=30mm,font=\scriptsize]
        {тип, число точек,\\число потомков,\\точки по индексу};
\end{tikzpicture}}

\vspace{0.5em}
{\small Один и тот же объект \texttt{Layer} умеет и читать, и фильтровать,
и считать --- а \texttt{Geometry} умеет разбирать сам себя рекурсивно.}
\end{frame}

% ======================= 11. ИЕРАРХИЯ ГЕОМЕТРИИ =======================
\begin{frame}{Иерархия геометрий: от контейнера к точке}
\footnotesize
\begin{columns}[T,onlytextwidth]
  \begin{column}{0.46\textwidth}
    \centering
    \begin{tikzpicture}[font=\scriptsize,
      lvl/.style={draw,rounded corners=2pt,minimum height=7mm,align=center,inner sep=3pt},
      ar/.style={-{Stealth[length=1.8mm]},draw=gray!70}]
      \node[lvl,fill=blue!10,text width=38mm] (a) {\texttt{POLYGON} / \texttt{MULTIPOLYGON}};
      \node[lvl,fill=blue!7,below=8mm of a,text width=38mm] (b) {\texttt{POLYGON} (кольцо)};
      \node[lvl,fill=blue!5,below=8mm of b,text width=38mm] (c) {\texttt{LINEARRING}};
      \node[lvl,fill=orange!10,below=8mm of c,text width=38mm] (d) {\texttt{POINT} $-$ координата};
      \draw[ar] (a)--(b); \draw[ar] (b)--(c); \draw[ar] (c)--(d);
      \node[right=4mm of a,anchor=west,font=\scriptsize] {уровень 0};
      \node[right=4mm of b,anchor=west,font=\scriptsize] {уровень 1};
      \node[right=4mm of c,anchor=west,font=\scriptsize] {уровень 2};
      \node[right=4mm of d,anchor=west,font=\scriptsize] {уровень 3};
    \end{tikzpicture}
  \end{column}
  \begin{column}{0.52\textwidth}
    \captionof{table}{\scriptsize Два метода --- это и есть рекурсия}
    \begin{tabular}{@{}lp{0.58\linewidth}@{}}
      \toprule
      Вызов & Что возвращает \\
      \midrule
      \texttt{GetGeometryCount()} & число непосредственных \emph{детей} \\
      \texttt{GetGeometryRef(i)} & $i$-й ребёнок \\
      \texttt{GetPointCount()}   & число точек \emph{внутри} узла \\
      \texttt{GetPoint(i)}       & координаты $i$-й точки \\
      \bottomrule
    \end{tabular}
    \vspace{0.4em}
    \begin{alertblock}{Ловушка, на которой спотыкаются все}\small
      Для \texttt{POLYGON} метод \texttt{GetPointCount()} возвращает
      \textbf{0}: точки лежат не в самом полигоне, а в его кольцах.
      Поэтому «просто пройтись по точкам» не работает --- нужен спуск
      по уровням.
    \end{alertblock}
  \end{column}
\end{columns}
\end{frame}

% ======================= 12. РАЗДЕЛ II =======================
\begin{frame}
  \begin{center}
    {\LARGE\bfseries Часть II. Чтение данных через OGR}\\[0.5em]
    {\large Пять программ --- пять шагов вглубь}\\[1.5em]
    {\small \textcolor{gray}{~55 минут}}
  \end{center}
\end{frame}

% ======================= 13. EXAMPLE01 КОД =======================
\begin{frame}[fragile]{$\texttt{example01.py}$: открыть и перечислить}
\begin{Verbatim}
import osgeo.ogr

shapefile = osgeo.ogr.Open("tl_2014_us_state.shp")
numLayers = shapefile.GetLayerCount()

print("Shapefile contains {} layers".format(numLayers))
print()

for layerNum in range(numLayers):
    layer = shapefile.GetLayer(layerNum)
    spatialRef = layer.GetSpatialRef().ExportToProj4()
    numFeatures = layer.GetFeatureCount()
    print("Layer {} has spatial reference {}".format(layerNum, spatialRef))
    print("Layer {} has {} features".format(layerNum, numFeatures))
    print()

    for featureNum in range(numFeatures):
        feature = layer.GetFeature(featureNum)
        featureName = feature.GetField("NAME")
        print("Feature {} has name {}".format(featureNum, featureName))
\end{Verbatim}
\end{frame}

% ======================= 14. EXAMPLE01 ВЫВОД =======================
\begin{frame}[fragile]{Что эта программа выводит на наших данных}
\begin{Verbatim}[numbers=none,fontsize=\scriptsize]
Shapefile contains 1 layers

Layer 0 has spatial reference +proj=longlat +datum=NAD83 +no_defs
Layer 0 has 56 features

Feature 0 has name West Virginia
Feature 1 has name Florida
Feature 2 has name Illinois
Feature 3 has name Minnesota
Feature 4 has name Maryland
Feature 5 has name Rhode Island
Feature 6 has name Idaho
Feature 7 has name New Hampshire
Feature 8 has name North Carolina
Feature 9 has name Vermont
Feature 10 has name Connecticut
Feature 11 has name Delaware
Feature 12 has name New Mexico
Feature 13 has name California
Feature 14 has name New Jersey
Feature 15 has name Wisconsin
   ...
Feature 55 has name Arizona
\end{Verbatim}
\vspace{0.2em}
{\footnotesize Три наблюдения: \textbf{(1)} в наборе \num{56} объектов ---
50 штатов, DC, Puerto Rico, Guam, American Samoa, Northern Mariana
Islands и US Virgin Islands; \textbf{(2)} порядок \emph{не} алфавитный;
\textbf{(3)} \texttt{feature 12} --- это \textbf{New Mexico},
\texttt{feature 55} --- \textbf{Arizona}.}
\end{frame}

% ======================= 15. РАЗБОР EXAMPLE01 =======================
\begin{frame}{Разбор по строкам: что здесь важно}
\footnotesize
\captionof{table}{\scriptsize Разбор ключевых строк \texttt{example01.py}}
\begin{tabular}{@{}p{0.34\linewidth}p{0.62\linewidth}@{}}
  \toprule
  Строка & Что происходит \\
  \midrule
  \texttt{ogr.Open("...shp")} &
    Возвращает \texttt{DataSource}. Чтение ленивое: файл открыт, но данные
    ещё не прочитаны. \\
  \texttt{GetLayerCount()} &
    У шейп-файла всегда \textbf{один} слой с именем
    \texttt{tl\_2014\_us\_state}. \\
  \texttt{GetLayer(0)} &
    Слой --- это «таблица»: строки-объекты и общая СК с геометрией. \\
  \texttt{layer.GetSpatialRef()} &
    СК слоя. Значение \texttt{None} возможно, только если СК не задана;
    здесь она задана. \\
  \texttt{GetFeatureCount()} &
    \num{56}. Это число, а не признак «данные уже в памяти». \\
  \texttt{GetFeature(12)} &
    Читает 13-ю запись: атрибуты \emph{и} ссылку на геометрию. \\
  \texttt{GetField("NAME")} &
    Достаёт поле по имени. В DBF имена чувствительны к регистру. \\
  \bottomrule
\end{tabular}
\vspace{0.5em}
{\scriptsize \textbf{Три способа пройти по объектам:} по индексу
(\texttt{GetFeature(i)}), «как по словарю» (\texttt{for f in layer})
или SQL-фильтром (\texttt{GetLayerByName(...).ExecuteSQL(...)}).}
\end{frame}

% ======================= 16. ПОРЯДОК ОБЪЕКТОВ =======================
\begin{frame}[fragile]{Порядок объектов: FIPS, а не алфавит}
\footnotesize
\begin{columns}[T,onlytextwidth]
  \begin{column}{0.5\textwidth}
    \captionof{table}{\scriptsize Первые объекты файла и их коды FIPS}
    \begin{tabular}{@{}rllc@{}}
      \toprule
      № & NAME & STATEFP & \\
      \midrule
      0  & West Virginia  & 54 & \\
      1  & Florida        & 12 & \\
      2  & Illinois       & 17 & \\
      3  & Minnesota      & 27 & \\
      4  & Maryland       & 24 & \\
      5  & Rhode Island   & 44 & \\
      6  & Idaho          & 16 & \\
      \dots & \dots & \dots & \\
      55 & Arizona        & 04 & \\
      \bottomrule
    \end{tabular}
    \vspace{0.4em}
    {\scriptsize Ни алфавит, ни сортировка по FIPS: это порядок,
     зафиксированный при сборке набора.}
  \end{column}
  \begin{column}{0.48\textwidth}
    \begin{block}{Практический вывод}\small
      \begin{itemize}\itemsep1pt
        \item «Нью-Мэксико --- это 12-й объект» верно только
              \emph{для этого файла};
        \item чтобы найти нужный штат, ищите по полю
              \texttt{STATEFP} / \texttt{STUSPS}, а не по номеру;
        \item в \texttt{NAME} есть и \texttt{Commonwealth of the Northern
              Mariana Islands} --- сравнивать строки надо аккуратно.
      \end{itemize}
    \end{block}
    \vspace{0.3em}
    \begin{Verbatim}[numbers=none,fontsize=\scriptsize]
lyr.SetAttributeFilter("NAME = 'Arizona'")
for f in lyr:
    print(f.GetFID())
\end{Verbatim}
  \end{column}
\end{columns}
\end{frame}

% ======================= 17. EXAMPLE02 =======================
\begin{frame}[fragile]{$\texttt{example02.py}$: атрибуты одного объекта}
\begin{Verbatim}
import osgeo.ogr

shapefile = osgeo.ogr.Open("tl_2014_us_state.shp")
layer = shapefile.GetLayer(0)
feature = layer.GetFeature(12)

print("Feature 12 has the following attributes:")
print()

attributes = feature.items()

for key,value in attributes.items():
    print(" {} = {}".format(key, value))

geometry = feature.GetGeometryRef()
geometryName = geometry.GetGeometryName()

print()
print("Feature's geometry data consists of a {}".format(geometryName))
\end{Verbatim}
\end{frame}

% ======================= 18. EXAMPLE02 ВЫВОД =======================
\begin{frame}[fragile]{Вывод: 14 полей и одна геометрия}
\begin{Verbatim}[numbers=none,fontsize=\scriptsize]
Feature 12 has the following attributes:

 REGION = 4
 DIVISION = 8
 STATEFP = 35
 STATENS = 00897535
 GEOID = 35
 STUSPS = NM
 NAME = New Mexico
 LSAD = 00
 MTFCC = G4000
 FUNCSTAT = A
 ALAND = 314161410324
 AWATER = 755712514
 INTPTLAT = +34.4346843
 INTPTLON = -106.1316181

Feature's geometry data consists of a POLYGON
\end{Verbatim}
\vspace{0.2em}
{\footnotesize \texttt{items()} --- «словарь всех полей»; порядок полей
в нём совпадает с порядком полей в DBF, то есть с порядком объявления
схемы, а не с алфавитным.}
\end{frame}

% ======================= 19. EXAMPLE02 РАЗБОР =======================
\begin{frame}[fragile]{Три способа достать атрибуты}
\footnotesize
\begin{columns}[T,onlytextwidth]
  \begin{column}{0.5\textwidth}
    \captionof{table}{\scriptsize API для атрибутов}
    \begin{tabular}{@{}p{0.56\linewidth}p{0.4\linewidth}@{}}
      \toprule
      Вызов & Результат \\
      \midrule
      \texttt{GetField("NAME")}      & \texttt{"New Mexico"} \\
      \texttt{GetField(6)}          & \texttt{"New Mexico"} \\
      \texttt{GetFieldAsString(6)}  & строка в любом случае \\
      \texttt{GetFieldAsInteger(6)} & целое, если поле числовое \\
      \texttt{items()}              & словарь всех полей \\
      \texttt{keys()}                & список имён полей \\
      \bottomrule
    \end{tabular}
  \end{column}
  \begin{column}{0.48\textwidth}
    \begin{Verbatim}[fontsize=\scriptsize]
geom = f.GetGeometryRef()
print(geom.GetGeometryName())
print(geom.GetEnvelope())
print(geom.Centroid().ExportToWkt())
\end{Verbatim}
    \vspace{0.2em}
    {\scriptsize \texttt{GetGeometryRef()} возвращает \emph{ссылку},
    а не копию: присваивание \texttt{g = f.GetGeometryRef()} не удваивает
    память.}
    \vspace{0.3em}
    {\scriptsize Для Нью-Мэксико \texttt{GetEnvelope()} даёт
     \texttt{(-109.050173, -103.001964,}\\
     \texttt{31.332172, 37.000293)} --- то есть
     \texttt{(xmin, xmax, ymin, ymax)} в градусах.}
  \end{column}
\end{columns}
\end{frame}

% ======================= 20. EXAMPLE03 =======================
\begin{frame}[fragile]{$\texttt{example03.py}$: рекурсивный разбор геометрии}
\begin{Verbatim}
import osgeo.ogr

def analyzeGeometry(geometry, indent=0):
    s = []
    s.append("  " * indent)
    s.append(geometry.GetGeometryName())
    if geometry.GetPointCount() > 0:
        s.append(" with %d data points" % geometry.GetPointCount())
    if geometry.GetGeometryCount() > 0:
        s.append(" containing:")

    print("".join(s))

    for i in range(geometry.GetGeometryCount()):
        analyzeGeometry(geometry.GetGeometryRef(i), indent+1)

shapefile = osgeo.ogr.Open("tl_2014_us_state.shp")
layer = shapefile.GetLayer(0)
feature = layer.GetFeature(12)
geometry = feature.GetGeometryRef()

analyzeGeometry(geometry)
\end{Verbatim}
\end{frame}

% ======================= 21. EXAMPLE03 ВЫВОД =======================
\begin{frame}[fragile]{Вывод: два уровня вложенности}
\begin{Verbatim}[numbers=none]
POLYGON containing:
  LINEARRING with 7554 data points
\end{Verbatim}

\vspace{0.3em}
{\small Отступ \texttt{"  " * indent} --- это и есть «дерево»,
напечатанное текстом. Разбор занял \textbf{19 строк} кода, и это
\textbf{единственный} алгоритм, который одинаково работает и для одного
кольца, и для Аляски с 50 островами.}

\vspace{0.5em}
\begin{columns}[T,onlytextwidth]
  \begin{column}{0.5\textwidth}
    \begin{block}{Почему рекурсия}\small
      Структура OGR --- это \emph{дерево}, а глубина зависит от типа
      объекта. Фиксированное число уровней здесь --- ловушка:
      у \texttt{MULTIPOLYGON} их четыре.
    \end{block}
  \end{column}
  \begin{column}{0.48\textwidth}
    \begin{Verbatim}[fontsize=\tiny]
MULTIPOLYGON containing:   # Alaska
  POLYGON containing:
    LINEARRING with 180 data points
  POLYGON containing:
    LINEARRING with 86 data points
  ... ещё 48 полигонов
\end{Verbatim}
  \end{column}
\end{columns}
\end{frame}

% ======================= 22. КОЛЬЦО =======================
\begin{frame}{Кольцо из 7554 точек}
\centering
\usePic{nm}\\[0.1cm]
\footnotesize
Нью-Мэксико --- почти идеальный прямоугольник в географических
координатах: границы идут по линиям $\varphi=\mathrm{const}$ и
$\lambda=\mathrm{const}$.
Синим пунктиром --- результат \texttt{GetExtent()}, оранжевые точки ---
выборка вершин из \texttt{GetPoint(i)}.
\end{frame}

% ======================= 23. СКОЛЬКО КОЛЕЦ =======================
\begin{frame}{Сколько колец у разных объектов}
\footnotesize
\captionof{table}{\scriptsize 12 объектов из 56 содержат больше одного кольца}
\begin{tabular}{@{}lrlrlr@{}}
  \toprule
  Объект & Колец & Вершин & Объект & Колец & Вершин \\
  \midrule
  Alaska                       & \num{50} & 14\,039 & American Samoa &  4 & 510 \\
  Hawaii                       & \num{17} & 3\,523  & New Jersey     &  3 & 6\,140 \\
  Commonwealth of the N. Marian Is. & \num{12} & 818 & New York  &  3 & 9\,976 \\
  California                   &  7 & 11\,131    & Puerto Rico    &  3 & 1\,565 \\
  Florida                      &  2 & 19\,798    & Kentucky       &  2 & 20\,306 \\
  Rhode Island                 &  2 & 2\,541     & U.S. Virgin Is. & 2 & 495 \\
  \bottomrule
\end{tabular}

\vspace{0.7em}
\begin{columns}[T,onlytextwidth]
  \begin{column}{0.5\textwidth}
    \begin{itemize}\itemsep1pt
      \item \textbf{Texas} --- рекордсмен по числу вершин: \num{58\,455};
      \item минимум --- \textbf{US Virgin Islands}: \num{495};
      \item \textbf{всего по набору} --- \num{855\,679} вершин.
    \end{itemize}
  \end{column}
  \begin{column}{0.48\textwidth}
    \begin{block}{Мораль}\small
      Код, который «знает», что у штата одно кольцо, сломается на Аляске.
      Рекурсия из \texttt{example03} написана так, чтобы работать с любым
      случаем --- это главный приём главы.
    \end{block}
  \end{column}
\end{columns}
\end{frame}

% ======================= 24. EXAMPLE04 =======================
\begin{frame}[fragile]{$\texttt{example04.py}$: рекурсивный поиск крайних точек}
\begin{Verbatim}[fontsize=\scriptsize]
import osgeo.ogr

def findPoints(geometry, results):
    for i in range(geometry.GetPointCount()):
        x,y,z = geometry.GetPoint(i)
        if results['north'] == None or results['north'][1] < y:
            results['north'] = (x,y)
        if results['south'] == None or results['south'][1] > y:
            results['south'] = (x,y)

    for i in range(geometry.GetGeometryCount()):
        findPoints(geometry.GetGeometryRef(i), results)

shapefile = osgeo.ogr.Open("tl_2014_us_state.shp")
layer = shapefile.GetLayer(0)
feature = layer.GetFeature(55)
geometry = feature.GetGeometryRef()

results = {'north' : None,
           'south' : None}

findPoints(geometry, results)

print("Northernmost point is ({:.4f}, {:.4f})".format(
      results['north'][0], results['north'][1]))
print("Southernmost point is ({:.4f}, {:.4f})".format(
      results['south'][0], results['south'][1]))
\end{Verbatim}
\end{frame}

% ======================= 25. EXAMPLE04 ВЫВОД =======================
\begin{frame}[fragile]{Вывод: Arizona, 9850 вершин}
\begin{Verbatim}[numbers=none]
Northernmost point is (-110.5007, 37.0043)
Southernmost point is (-111.0006, 31.3322)
\end{Verbatim}
\vspace{0.2em}
{\footnotesize Это \textbf{feature 55} --- Arizona. Одна функция ---
два результата.}
\vspace{0.2em}
\centering
\usePic{az}\\[0.1cm]
{\footnotesize Красный отрезок --- путь между найденными крайними
вершинами; он же совпадает с границами \texttt{GetEnvelope()}
по вертикали: 37.0043 и 31.3322.}
\end{frame}

% ======================= 26. АНАЛИЗ EXAMPLE04 =======================
\begin{frame}[fragile]{Что здесь сделано хорошо и что можно улучшить}
\footnotesize
\begin{columns}[T,onlytextwidth]
  \begin{column}{0.5\textwidth}
    \begin{block}{Хорошо}\small
      \begin{itemize}\itemsep1pt
        \item \textbf{одна функция --- два результата}: словарь
              \texttt{results} играет роль «накопителя»;
        \item \textbf{два независимых прохода}: сначала все точки текущего
              узла, потом рекурсия по детям;
        \item \textbf{сравнение через \texttt{None}}: алгоритм работает
              и с первого вызова.
      \end{itemize}
    \end{block}
  \end{column}
  \begin{column}{0.48\textwidth}
    \begin{alertblock}{Что улучшить}\small
      \begin{itemize}\itemsep1pt
        \item \texttt{== None} лучше писать как \texttt{is None};
        \item есть только «север» и «юг» --- запад и восток добавляются
              так же (сравнение по $x$);
        \item «север» и «юг» условны: речь о максимуме и минимуме широты
              \emph{в градусах};
        \item кольцо замкнуто и идёт по границе, поэтому экстремум по $y$
              можно искать вообще без Python-цикла.
      \end{itemize}
    \end{alertblock}
  \end{column}
\end{columns}

\vspace{0.5em}
\begin{columns}[T,onlytextwidth]
  \begin{column}{0.5\textwidth}
    \begin{block}{Тот же результат силами OGR}\small
      \texttt{geometry.GetEnvelope()} или \texttt{layer.GetExtent()} ---
      без единой вершины в Python-коде.
    \end{block}
  \end{column}
  \begin{column}{0.48\textwidth}
    \begin{Verbatim}[fontsize=\scriptsize]
g = lyr.GetFeature(55).GetGeometryRef()
xmin, xmax, ymin, ymax = g.GetEnvelope()
print(f"{ymin:.4f} {ymax:.4f}")
# 31.3322 37.0043  <- совпало
\end{Verbatim}
  \end{column}
\end{columns}
\end{frame}

% ======================= 27. ПРОИЗВОДИТЕЛЬНОСТЬ =======================
\begin{frame}{Производительность: во что обходится обход}
\footnotesize
Измерено на нашем файле (Python 3, GDAL 3.13, один поток):

\vspace{0.2em}
\captionof{table}{\scriptsize Замеры на feature 55 (Arizona, 9850 вершин)}
\begin{tabular}{@{}p{0.52\linewidth}rl@{}}
  \toprule
  Операция & Время & Комментарий \\
  \midrule
  \texttt{ogr.Open} + \texttt{GetLayerCount} & $\sim$0 мс &
    файл открыт, но не прочитан \\
  \texttt{GetFeature(55)} (первое обращение) & 11.7 мс &
    читается запись из \texttt{.shp} \\
  \texttt{GetPoint(i)} в цикле по Python & 5.0 мс &
    9850 вызовов через SWIG \\
  \texttt{GetPoints()} целиком (OGR) & 0.9 мс &
    \textbf{в 5.6 раза быстрее} \\
  \texttt{GetEnvelope()} & 0.02 мс &
    \textbf{в 250 раз быстрее} \\
  \midrule
  рекурсивный обход всех 56 объектов & 15 мс &
    все \num{855\,679} вершин \\
  \bottomrule
\end{tabular}

\vspace{0.7em}
\begin{alertblock}{Правило}\small
  Простой \texttt{GetPoint(i)} удобен для обучения и точечных запросов.
  Для массовых операций спрашивайте у OGR сразу всё
  (\texttt{GetPoints()}, \texttt{GetEnvelope()}, \texttt{ExportToWkb()}),
  а не перебирайте вершины по одной.
\end{alertblock}
\end{frame}

% ======================= 28. ПОДВОДНЫЕ КАМНИ =======================
\begin{frame}[fragile]{Подводные камни, о которые стоит споткнуться}
\footnotesize
\begin{columns}[T,onlytextwidth]
  \begin{column}{0.52\textwidth}
    \begin{enumerate}\itemsep1pt
      \item \textbf{\texttt{GetPointCount()} у \texttt{POLYGON} = 0.}
            Точки лежат в кольцах. Проверили:
            \texttt{(0, 1, LINEARRING, 7554)}.
      \item \textbf{Третья координата.} \texttt{GetPoint(i)} возвращает
            \texttt{(x, y, z)}, а в нашем файле $z = 0$. Распаковка
            \texttt{x, y, z = ...} работает, \texttt{x, y = ...} --- нет.
      \item \textbf{Нет обработки ошибок.} \texttt{ogr.Open} вернёт
            \texttt{None}, и через три строки вы получите
            \texttt{AttributeError} вместо внятного сообщения.
      \item \textbf{Устаревший API.} Современный GDAL (3.7+) по умолчанию
            бросает исключения вместо \texttt{None} --- включите
            \texttt{gdal.UseExceptions()} явно.
    \end{enumerate}
  \end{column}
  \begin{column}{0.46\textwidth}
    \begin{alertblock}{И ещё}\small
      \begin{itemize}\itemsep1pt
        \item \textbf{порядок} объектов не гарантирован ничем;
        \item \textbf{кольцо замкнуто}: первая и последняя точки совпадают
              (проверили: \texttt{closed=True});
        \item \textbf{антимеридиан}: extent набора уходит в
              $\lambda\in[-179.23;\,179.86]$ --- Аляска «перескакивает»
              на $+180^\circ$, и наивная проекция даст артефакты;
        \item \textbf{доминантное кольцо}: в многоугольнике с «дырами»
              первое кольцо --- внешнее, остальные --- внутренние.
      \end{itemize}
    \end{alertblock}
    \vspace{0.2em}
    \begin{Verbatim}[fontsize=\tiny,numbers=none]
r = g.GetGeometryRef(0)
n = r.GetPointCount()
print(r.GetPoint(0) == r.GetPoint(n-1))
# True
\end{Verbatim}
  \end{column}
\end{columns}
\end{frame}

% ======================= 29. РАЗДЕЛ III =======================
\begin{frame}
  \begin{center}
    {\LARGE\bfseries Часть III. Геодезия}\\[0.5em]
    {\large Откуда в векторной главе вдруг появилась тригонометрия}\\[1.5em]
    {\small \textcolor{gray}{~30 минут}}
  \end{center}
\end{frame}

% ======================= 30. ЗАЧЕМ РАССТОЯНИЕ =======================
\begin{frame}{Зачем вдруг расстояние между точками?}
\footnotesize
\begin{columns}[T,onlytextwidth]
  \begin{column}{0.52\textwidth}
    \begin{itemize}\itemsep2pt
      \item Данные из шейп-файла лежат в \textbf{градусах}. Чтобы сравнить
            «длину границы» с «площадью» или ответить на вопрос
            «насколько далеко?», градусы надо превратить в метры.
      \item Геодезия решает это на \textbf{эллипсоиде} --- но сложно.
      \item Простейшая математическая модель Земли --- \textbf{сфера},
            одна константа: радиус $R$.
      \item Путь по поверхности сферы, соединяющий две точки, --- дуга
            \textbf{большого круга}. Она кратчайшая.
      \item Значит, задача сводится к поиску длины этой дуги.
    \end{itemize}
  \end{column}
  \begin{column}{0.46\textwidth}
    \begin{block}{Две координаты из примера}\small
      \begin{tabular}{@{}lll@{}}
        \toprule
         & широта & долгота \\
        \midrule
        точка 1 & 42.0095 & $-122.3782$ \\
        точка 2 & 32.5288 & $-117.2049$ \\
        \bottomrule
      \end{tabular}
      \medskip

      Это [\emph{Crater Lake}, Орегон] и [\emph{San Diego}, Калифорния].
      Наклон $\Delta\varphi = -9.48^\circ$, разброс по долготе
      $\Delta\lambda = +5.17^\circ$.
    \end{block}
  \end{column}
\end{columns}
\end{frame}

% ======================= 31. EXAMPLE05 КОД =======================
\begin{frame}[fragile]{$\texttt{example05.py}$: гаверсинус в 20 строк}
\begin{Verbatim}
import math

lat1 = 42.0095
long1 = -122.3782

lat2 = 32.5288
long2 = -117.2049

rLat1 = math.radians(lat1)
rLong1 = math.radians(long1)

rLat2 = math.radians(lat2)
rLong2 = math.radians(long2)

dLat = rLat2 - rLat1
dLong = rLong2 - rLong1
a = math.sin(dLat/2)**2 + math.cos(rLat1) * math.cos(rLat2) \
                        * math.sin(dLong/2)**2
c = 2 * math.atan2(math.sqrt(a), math.sqrt(1-a))
distance = 6371 * c

print("Great circle distance is {:.0f} kilometers".format(distance))
\end{Verbatim}
\end{frame}

% ======================= 32. РАЗБОР КОДА =======================
\begin{frame}{Разбор формулы по шагам}
\scriptsize
\begin{columns}[T,onlytextwidth]
  \begin{column}{0.56\textwidth}
    \captionof{table}{\scriptsize Что происходит в каждой строке}
    \setlength{\tabcolsep}{3pt}
    \begin{tabular}{@{}p{0.36\linewidth}p{0.60\linewidth}@{}}
      \toprule
      Строка & Смысл \\
      \midrule
      \texttt{math.radians(...)} &
        градусы $\to$ радианы: тригонометрия Python работает только
        в радианах \\
      \texttt{dLat}, \texttt{dLong} &
        разности широт и долгот в радианах: $-0.165469$ и $0.090291$ \\
      \texttt{math.sin(dLat/2)**2} &
        $\mathrm{hav}\,\Delta\varphi$ --- «полухорд» \\
      \texttt{math.cos(rLat1)*\dots} &
        $\cos\varphi_1\cos\varphi_2$ --- поправка на широту \\
      \texttt{a} &
        значение подкоренного выражения: \num{0.0081053792} \\
      \texttt{atan2(\dots)} &
        устойчивый «обратный синус»: не теряет точность при $a\approx1$ \\
      \texttt{6371 * c} &
        средний радиус Земли $\times$ центральный угол \\
      \bottomrule
    \end{tabular}
  \end{column}
  \begin{column}{0.42\textwidth}
    \begin{block}{Зачем \texttt{atan2}, а не \texttt{asin}}\small
      Наивный вариант
      \texttt{c = 2*math.asin(math.sqrt(a))} ломается, когда точки близки
      или совпадают: $\sqrt a\to1$, производная \texttt{asin} уходит
      в бесконечность. \texttt{atan2} считает ту же величину устойчиво.
    \end{block}
    \vspace{0.3em}
    \begin{block}{Про радиус}\small
      \texttt{6371} км --- средний радиус по Объединённому геодезическому
      эллипсоиду. От его выбора напрямую зависит точность.
    \end{block}
  \end{column}
\end{columns}
\end{frame}

% ======================= 33. ФОРМУЛА =======================
\begin{frame}{Формула целиком}
\begin{equation*}
  \mathrm{hav}\,\sigma = \sin^2\frac{\sigma}{2}
  = \mathrm{hav}\,\Delta\varphi
  + \cos\varphi_1 \cdot \cos\varphi_2 \cdot \mathrm{hav}\,\Delta\lambda
\end{equation*}
\begin{equation*}
  \sigma = 2\,\mathrm{atan2}\!\left(\sqrt{a},\ \sqrt{1-a}\right),
  \qquad d = R\,\sigma
\end{equation*}

\vspace{0.4em}
\footnotesize
\begin{columns}[T,onlytextwidth]
  \begin{column}{0.5\textwidth}
    \begin{itemize}\itemsep1pt
      \item $\mathrm{hav}(\theta)\equiv\sin^2(\theta/2)$ --- «полухорд»;
      \item $\sigma$ --- \textbf{центральный угол} между точками
            (угол в центре сферы);
      \item $a$ --- значение в скобках; оно всегда в $[0;\,1]$,
            поэтому оба корня вещественны;
      \item $d$ --- длина дуги: \textbf{радиусы} умножаются на
            \textbf{угол в радианах}.
    \end{itemize}
  \end{column}
  \begin{column}{0.48\textwidth}
    \begin{block}{Почему $\cos\varphi_1\cos\varphi_2$}\small
      Разница долгот $\Delta\lambda$ сама по себе не даёт расстояния:
      на экваторе $1^\circ$ долготы $\approx$ 111 км, а на $60^\circ$
      с.ш. --- всего 55 км. Множитель $\cos\varphi_1\cos\varphi_2$ это
      учитывает.
      \medskip

      На наших данных:
      $\cos 42.0095^\circ\cdot\cos 32.5288^\circ = 0.626468$.
    \end{block}
  \end{column}
\end{columns}
\end{frame}

% ======================= 34. ГЕОМЕТРИЯ =======================
\begin{frame}{Что означает формула геометрически}
\centering
\usePic{gc}\\[0.1cm]
\footnotesize
Три «числа» в формуле --- это стороны прямоугольного сферического
треугольника: $\Delta\varphi$ (разница широт), $\Delta\lambda$ (разница
меридианов) и $\sigma$ (центральный угол между точками).
\end{frame}

% ======================= 35. ВЫВОД ФОРМУЛЫ =======================
\begin{frame}{Откуда берётся формула: тригонометрия в 4 строки}
\footnotesize
На сфере радиуса $R$ длины дуг пропорциональны углам, поэтому достаточно
работать с \emph{углами} треугольника со сторонами
$\bigl(\tfrac{\pi}{2}-\varphi_1\bigr)$,
$\bigl(\tfrac{\pi}{2}-\varphi_2\bigr)$ и $\sigma$.

\begin{enumerate}\itemsep1pt
  \item \textbf{Сферическая теорема косинусов} для прямого угла при $A$:
        \[
          \cos\sigma = \sin\varphi_1\sin\varphi_2
          + \cos\varphi_1\cos\varphi_2\cos\Delta\lambda .
        \]
  \item Подставляем $\cos\Delta\lambda = 1 - 2\,\mathrm{hav}\,\Delta\lambda$
        и $\cos\sigma = 1 - 2\,\mathrm{hav}\,\sigma$:
        \[
          1 - 2\,\mathrm{hav}\,\sigma
          = 1 - 2\,\mathrm{hav}\,\Delta\varphi
            - 2\cos\varphi_1\cos\varphi_2\,\mathrm{hav}\,\Delta\lambda .
        \]
  \item Сокращаем единицы и делим на $2$:
        \[
          \mathrm{hav}\,\sigma = \mathrm{hav}\,\Delta\varphi
          + \cos\varphi_1\cos\varphi_2\,\mathrm{hav}\,\Delta\lambda .
        \]
  \item Берём $\sigma = 2\,\mathrm{atan2}(\sqrt a,\sqrt{1-a})$
        и умножаем на $R$.
\end{enumerate}

\vspace{0.3em}
\begin{block}{Откуда взялся именно «полухорд»}\small
  Тождества $\cos\Delta\lambda = 1-2\sin^2\frac{\Delta\lambda}{2}$ и
  $\cos\sigma = 1-2\sin^2\frac{\sigma}{2}$ --- это формулы половинного
  угла. Именно они убирают из формулы \emph{радианные разности}
  $\Delta\varphi$, $\Delta\lambda$, которые иначе пришлось бы считать
  через \texttt{radians}. Это делает формулу устойчивой: $\mathrm{hav}$
  всегда остаётся маленьким числом.
\end{block}
\end{frame}

% ======================= 36. ЧИСЛА =======================
\begin{frame}{Числа примера по шагам}
\footnotesize
\captionof{table}{\scriptsize Промежуточные значения (вычислены независимо, совпадают с программой)}
\begin{tabular}{@{}lrl@{\qquad}lrl@{}}
  \toprule
  Величина & & Значение & Величина & & Значение \\
  \midrule
  $\Delta\varphi$ & $=$ & $-0.165469$ рад & $\cos\varphi_1\cos\varphi_2$ & $=$ & $0.626468$ \\
  $\Delta\varphi$ & $=$ & $-9.4807^\circ$ & $\mathrm{hav}\,\Delta\varphi$ & $=$ & $0.0068294$ \\
  $\Delta\lambda$ & $=$ & $0.090291$ рад & $\mathrm{hav}\,\Delta\lambda$ & $=$ & $0.0020367$ \\
  $\Delta\lambda$ & $=$ & $5.1733^\circ$ & $a$ & $=$ & $0.0081053792$ \\
  $\sigma$ & $=$ & $0.1803039$ рад & $\sigma$ & $=$ & $10.3307^\circ$ \\
  $R$ & $=$ & 6371 км & $d = R\sigma$ & $=$ & \textbf{1148.716 км} \\
  \bottomrule
\end{tabular}

\vspace{0.8em}
{\small Программа печатает с форматом \texttt{\{:.0f\}} --- то есть
\num{1149} kilometers. Округление до целого скрывает всю точность формулы:
полезно знать и то, и другое.}

\vspace{0.5em}
\begin{block}{Для самопроверки}\small
  Подставьте в код $R = 6378.137$: получится \num{1150.003} км.
  Разница \num{1.29} км целиком объясняется выбором радиуса.
\end{block}
\end{frame}

% ======================= 37. ПРОВЕРКА =======================
\begin{frame}{Проверка на настоящей геодезии}
\centering
\usePic{dc}\\[0.1cm]
\footnotesize
Эталон --- решение на эллипсоиде WGS84 (\texttt{pyproj}, метод геодезических
линий): \num{1147.3} км. Гаверсинус со средним радиусом отличается на
\num{+0.12\,\%} --- это примерно 1.4 км на дистанции в 1150 км, то есть
на дальних расстояниях приемлемо.
\end{frame}

% ======================= 38. ТАБЛИЦА РЕКОМЕНДАЦИЙ =======================
\begin{frame}{Когда какой метод уместен}
\footnotesize
\captionof{table}{\scriptsize Сравнение подходов на нашем примере}
\renewcommand{\arraystretch}{0.95}
\begin{tabular}{@{}p{0.46\linewidth}rl@{}}
  \toprule
  Метод & Результат & Ошибка \\
  \midrule
  Геодезия на эллипсоиде (PROJ/WGS84) & 1147.3 км & --- (эталон) \\
  Плоская СК UTM, зона 11N & 1147.8 км & $+0.04\,\%$ \\
  \textbf{Гаверсинус, $R=6371$} & \textbf{1148.7 км} & $+0.12\,\%$ \\
  Гаверсинус, $R=6378.137$ & 1150.0 км & $+0.24\,\%$ \\
  «111.32 км на градус» с $\cos\varphi$ & 1150.6 км & $+0.29\,\%$ \\
  Евклидово в градусах $\times(6371\pi/180)$ & 1200.9 км & $+4.67\,\%$ \\
  \bottomrule
\end{tabular}

\vspace{0.2em}
\begin{columns}[T,onlytextwidth]
  \begin{column}{0.5\textwidth}
    \begin{block}{Берите}\small
      \begin{itemize}\itemsep1pt
        \item \textbf{гаверсинус} --- почти всегда достаточно: быстро
              и без внешних зависимостей, точность $\sim0.5\,\%$
              на любой дистанции;
        \item \textbf{UTM} --- когда нужны метры в плоскости;
        \item \textbf{геодезию PROJ} --- когда важны сантиметры.
      \end{itemize}
    \end{block}
  \end{column}
  \begin{column}{0.48\textwidth}
    \begin{alertblock}{Осторожно}\small
      \begin{itemize}\itemsep1pt
        \item последний способ --- ловушка: «градусы $\times$ 111»
              даёт ошибку $4.7\,\%$;
        \item при \texttt{atan2} не забывайте про \textbf{радианы};
        \item UTM-зона --- \emph{полоса}, а не глобальная проекция.
      \end{itemize}
    \end{alertblock}
  \end{column}
\end{columns}
\end{frame}

% ======================= 39. ТИПИЧНЫЕ ОШИБКИ =======================
\begin{frame}[fragile]{Чек-лист «как не ошибиться»}
\footnotesize
\begin{columns}[T,onlytextwidth]
  \begin{column}{0.5\textwidth}
    \begin{enumerate}\itemsep1pt
      \item \textbf{Градусы вместо радиан.}
            \texttt{math.sin(9.48)} вместо \texttt{math.radians(9.48)} ---
            тихая ошибка на 3--4 порядка.
      \item \textbf{Перепутаны широта и долгота.} В PROJ и OGR порядок
            \textbf{(x, y) = (долгота, широта)}; в \texttt{GetPoint} ---
            тоже. В списке людей --- наоборот.
      \item \textbf{Радиус в километрах, результат в метрах}
            (или наоборот).
      \item \textbf{Нет проверки на \texttt{None}} после
            \texttt{ogr.Open} и \texttt{GetFeature}.
      \item \textbf{Забыли про рекурсию}: у \texttt{POLYGON} точки
            напрямую не читаются.
    \end{enumerate}
  \end{column}
  \begin{column}{0.48\textwidth}
    \begin{Verbatim}[fontsize=\scriptsize]
# метры -> километры явно
d_km = 6371.0 * c

# проверка "на глаз"
assert 0.0 <= a <= 1.0
assert d_km > 0

# безопасное чтение
ds = ogr.Open(path)
if ds is None:
    raise IOError(path)
f = ds.GetLayer(0).GetFeature(55)
if f is None:
    raise KeyError("нет объекта 55")
\end{Verbatim}
    \vspace{0.2em}
    \begin{block}{Хорошая привычка}\small
      Перед расчётом расстояния напечатайте \texttt{ExportToProj4()}.
      Если там не градусы --- формула гаверсинуса бессмысленна.
    \end{block}
  \end{column}
\end{columns}
\end{frame}

% ======================= 40. ДАЛЬШЕ =======================
\begin{frame}[fragile]{Куда движется глава 3}
\footnotesize
\begin{columns}[T,onlytextwidth]
  \begin{column}{0.5\textwidth}
    \begin{block}{Что мы сделали «руками»}\small
      \begin{itemize}\itemsep1pt
        \item чтение шейп-файла --- \texttt{osgeo.ogr} (OGR, C++);
        \item разбор геометрии --- те же рекурсивные методы OGR;
        \item расстояние --- своя формула на \texttt{math}.
      \end{itemize}
    \end{block}
    \vspace{0.2em}
    \begin{alertblock}{Это и делают инструменты главы 3}\small
      \begin{itemize}\itemsep1pt
        \item \texttt{Fiona} --- векторные данные (обёртка над OGR);
        \item \texttt{Pyproj} --- СК и геодезические расчёты;
        \item \texttt{Shapely} --- геометрия и операции над ней;
        \item \texttt{GDAL} --- растры.
      \end{itemize}
    \end{alertblock}
  \end{column}
  \begin{column}{0.48\textwidth}
    \begin{Verbatim}[fontsize=\scriptsize]
# example05, но на эллипсоиде
from pyproj import Geod

g = Geod(ellps="WGS84")
_, _, d = g.inv(-122.3782, 42.0095,
                -117.2049, 32.5288)
print(d / 1000)
# 1147.303...  <- точнее

# и на Shapely: та же сфера
from shapely.geometry import Point
p1 = Point(-122.3782, 42.0095)
p2 = Point(-117.2049, 32.5288)
print(p1.distance(p2) / 1000)
# 1148.7...
\end{Verbatim}
    \vspace{0.3em}
    {\scriptsize То есть та же математика, только реализованная один раз
     и со всеми проверками.}
  \end{column}
\end{columns}
\end{frame}

% ======================= 41. ИТОГИ =======================
\begin{frame}{Итоги: пять тезисов}
\footnotesize
\begin{enumerate}\itemsep1pt
  \item \textbf{Шейп-файл} --- это пять файлов (геометрия, индекс,
        атрибуты, СК), и порядок объектов в нём произвольный: ищите
        объекты по полям, а не по номерам.
  \item \textbf{OGR} даёт стройную объектную модель
        \texttt{DataSource $\to$ Layer $\to$ Feature $\to$ Geometry};
        все четыре уровня нужны для простой задачи.
  \item \textbf{Геометрия --- это дерево}, а не список точек: рекурсия
        по \texttt{GetGeometryCount()}/\texttt{GetGeometryRef()}
        обязательна, а \texttt{GetPointCount()} у \texttt{POLYGON}
        равен нулю.
  \item \textbf{Рекурсия решает и «бытовые» задачи}: крайние точки штата
        находятся за 20 строк --- а \texttt{GetEnvelope()} даёт то же
        самое за одну.
  \item \textbf{Гаверсинус} --- переход от градусов к километрам через
        сферу: ошибка $\sim0.5\,\%$ вместо $4.7\,\%$ у «градусных»
        формул. Для сантиметров нужен эллипсоид PROJ.
\end{enumerate}

\vspace{0.6em}
\begin{block}{Ключевая мысль главы}\small
  Геопространственные данные --- это \emph{структура}, а не файл.
  Пока вы понимаете, как устроен объект изнутри (поля + дерево геометрии
  + СК), любой формат читается одинаково.
\end{block}
\end{frame}

% ======================= 42. ДОМАШНЕЕ ЗАДАНИЕ =======================
\begin{frame}[fragile]{Практика}
\footnotesize
\textbf{1. Запустить всё.} Положить рядом \texttt{tl\_2014\_us\_state.shp}
и выполнить пять примеров. Сверить вывод с числами этой презентации
(56 объектов, feature 12 = New Mexico, feature 55 = Arizona, 1149 км).

\textbf{2. Найти, а не угадать.} Модифицировать \texttt{example01.py}
так, чтобы он печатал только объекты с
\texttt{ALAND > 1\,000\,000\,000} м$^2$ (10 штатов-гигантов), и проверить,
что среди них есть Texas.

\textbf{3. Рекурсия.} Написать функцию, которая для заданного номера
объекта печатает его \emph{тип} и \emph{число колец}. Прогнать по всем
56 объектам: 44 объекта с одним кольцом и 12 --- с несколькими.

\textbf{4. Экстремумы.} Дополнить \texttt{example04.py} поиском самых
западной и восточной вершин и сравнить результат с
\texttt{GetEnvelope()}.

\textbf{5. Геодезия.} Сравнить три реализации расстояния между теми же
двумя точками: \texttt{example05.py}, \texttt{pyproj.Geod} и
\texttt{Shapely} (\texttt{Point.distance}). Построить таблицу отклонений;
проверить, что при \texttt{R = 6378.137} гаверсинус даёт 1150.0 км, а на
эллипсоиде --- 1147.3 км.

\textbf{6.} \textit{(дополнительно)} Посчитать площадь Нью-Мэксико двумя
способами: \texttt{geometry.GetArea()} в квадратных градусах и «через UTM»
в квадратных километрах, затем сравнить с \texttt{ALAND} из атрибутов.
\end{frame}

% ======================= 43. ВОПРОСЫ =======================
\begin{frame}
  \begin{center}
    {\Huge Вопросы?}\\[1.2em]
    {\large Пять файлов, один шейп-файл, \num{855\,679} вершин}\\[0.8em]
    {\normalsize
     \begin{tabular}{ll}
       \texttt{example01.py} & слои, СК, объекты \\
       \texttt{example02.py} & атрибуты и имя геометрии \\
       \texttt{example03.py} & дерево геометрии \\
       \texttt{example04.py} & крайние точки \\
       \texttt{example05.py} & расстояние по большому кругу
     \end{tabular}}
  \end{center}
\end{frame}

\end{document}